% ---------------------------------------------------------------------
\section{与 OPE 结果的一致性}\label{sec:disall-consistency}

OPE 的 Wilson 系数 $C^{(n)}_{a,i}$（\cref{ch:dis}）与动量空间证明中由部分子靶构造的系数函数（\cref{prop:disall-h5}）必须一致；一致的是作为分布的扭度 2 部分，而不是逐点的余项。

\begin{corollary}[与 OPE 系数的一致性 \tagR]\label{cor:disall-consistency}
（$R_{\mathrm{OPE}}$：以 OPE 为前提而严格。QCD 中含夸克的微扰 OPE 是全阶论证（\cref{thm:dis-ope}），严格证明只对标量理论与纯 Yang--Mills 的规范不变算符给出；对强子态还须非微扰 OPE（\cref{asm:intro-np-ope}）。）
以\cref{prop:pdf-partonic-dr}的构造定义动量空间的系数函数 $C^{\mathrm{mom}}_{a,i}$：在纯维数正规化中对在壳无质量部分子靶 $j$，$\hat F^{(j)}_a/z=\sum_iC^{\mathrm{mom}}_{a,i}\otimes f_{i/j}$，$f_{i/j}$ 为 $\MSbar$ 部分子 PDF。
则对 OPE 有定义的一切 $n$，$\int_0^1\dd z\,z^{n-1}C^{\mathrm{mom}}_{a,i}(z)=C^{(n)}_{a,i}$；因此 $C^{\mathrm{mom}}_{a,i}$ 就是\cref{thm:pdf-dis-xspace}中以其矩唯一确定的分布 $C_{a,i}$，
\cref{thm:disall-factorization}与\cref{thm:pdf-dis-xspace}的扭度 2 部分作为分布相同；逐点的余项只受全阶论证\cref{thm:disall-factorization}控制。（本推论只涉及部分子态的系数函数。）
\uses{thm:pdf-dis-xspace,prop:dis-wilson-partons,cor:appK-pdf,prop:pdf-partonic-dr,lem:prelim-mellin,thm:pdf-moments-local}
\end{corollary}

\begin{proof}
(1) 由 Mellin 卷积定理（\cref{lem:prelim-mellin}），$\int_0^1\dd z\,z^{n-1}(C\otimes f)(z)=C_nf_n$，$X_n\equiv\int_0^1\dd z\,z^{n-1}X(z)$。对定义式取矩：
$\int_0^1\dd z\,z^{n-2}\hat F^{(j)}_2=\sum_iC^{\mathrm{mom}}_{2,i,n}\,f_{i/j,n}$。
(2) OPE 一侧，\cref{prop:dis-wilson-partons}（以 OPE 为前提而严格）给出 $\int_0^1\dd z\,z^{n-2}\hat F^{(j)}_2=\sum_iC^{(n)}_{2,i}\langle j|O^{(n)}_i|j\rangle$，
其中在壳无质量部分子的算符矩阵元的全部圈修正都无标度，裸矩阵元等于 $\delta_{ij}$，重整化后等于 $\delta_{ij}$ 加纯 $\MSbar$ 极点。由\cref{thm:pdf-moments-local}，
重整化的扭度 2 定域算符矩阵元等于重整化 PDF 的矩；对部分子靶，二者都等于“裸值 $\delta_{ij}$ 乘以同一个 $Z^{-1}$”（\cref{prop:pdf-partonic-dr}），故 $\langle j|O^{(n)}_i|j\rangle=f_{i/j,n}$。
(3) 两式相减：$\sum_i\big(C^{\mathrm{mom}}_{2,i,n}-C^{(n)}_{2,i}\big)f_{i/j,n}=0$ 对一切部分子靶 $j$ 成立；矩阵 $(f_{i/j,n})=\mathbb 1+\order{\as}$ 逐阶可逆，故逐阶地 $C^{\mathrm{mom}}_{2,i,n}=C^{(n)}_{2,i}$。$F_L$（以及 $F_1$）同理。
(4) $C^{\mathrm{mom}}_{2,i}$ 是 $[0,1]$ 上的有限阶分布（$\delta(1-z)$、加号分布与可积函数）；由\cref{cor:appK-pdf}的论证，偶数阶矩（$n\ge N_1$）唯一地确定这样的分布，
所以 $C^{\mathrm{mom}}_{2,i}$ 等于\cref{thm:pdf-dis-xspace}中矩为 $C^{(n)}_{2,i}$ 的那个分布。于是两个定理中 $x$ 空间的主项逐项相同。
(5) 反过来，矩的一致不能推出余项逐点一致：对每个固定 $n$ 矩为 $\order{1/Q^2}$ 的函数逐点可以是 $\order{1}$（例如 $\sin(Q^2x)$；同类反例见\cref{prop:disall-two-routes}证明的 (c)）。
所以\cref{thm:pdf-dis-xspace}只控制分布意义下的余项，固定 $x$ 处的逐点余项是全阶论证\cref{thm:disall-factorization}的内容。
\end{proof}

\noindent
在 $\order{\as}$ 这一一致性已被逐项检验：在壳部分子前向 Compton 振幅硬区的 Taylor 系数，对 $\varepsilon$ 精确地等于割图结果的 Mellin 矩（夸克靶 $N\le8$，胶子靶偶数 $N\le8$），
重整化后给出\cref{ch:dis}的闭式（\cref{ex:dis-oneloop-ope}）；下一节 $x$ 空间系数函数的矩也与之一致。   % VF-DIS-13
共振区 $W^2\sim\Lambda^2$ 只能在对 $x$ 平滑后与扭度展开比较（Bloom--Gilman 对偶，猜想层次）。
